% !TeX encoding = UTF-8
% !TeX spellcheck = en_US
% Letter with HMGU template
% March 2016
% Based on letter template for ICube from Vincent Mazet
% Infos and bugs julien.pontabry@helmholtz-muenchen.de


% Load letter class (change the paper size if you need)
\documentclass[a4paper,10pt,twosided]{letter}

% English language is used in the letter
\usepackage[english]{babel}

% Set the date of the letter (use \today for today)
\date{\today}

% Layout file of the template
% !TeX encoding = UTF-8
% !TeX spellcheck = en_US
% Letter template for HMGU
% March 2016
% Based on letter template for ICube from Vincent Mazet
% Infos and bugs julien.pontabry@helmholtz-muenchen.de

% Packages
\usepackage[utf8]{inputenc} 
\usepackage[T1]{fontenc}
\usepackage{graphicx}
\usepackage{fancyhdr}
\usepackage{color}
\usepackage{ifthen}
\usepackage{helvet}

% Layout
\voffset -1in
\topmargin 10mm
\headheight 5mm
\headsep 35mm
\textheight 210mm
\footskip 20mm
\hoffset -1in
\oddsidemargin 15mm
\evensidemargin 15mm
\textwidth 135mm
\marginparsep -180mm
\marginparwidth 350mm
\parindent 0mm

% Font sans-serif
\renewcommand{\familydefault}{\sfdefault}

% Colors
\definecolor{grey}{RGB}{128,128,128}

% Logos
\newcommand{\hmguen}{\includegraphics[width=115mm]{hmgu_color_englisch.eps} \\[6ex]{\scriptsize\color{grey}Helmholtz Zentrum München - P.O. box 11 29 - 85758 Neuherberg}}

% Headers and footers
\pagestyle{fancy}
\renewcommand{\headrulewidth}{0pt}
\lhead{\hspace*{-10mm}\ifthenelse{\value{page}=1}{\hmguen}{}}
\chead{}
\lfoot{}
\cfoot{}
\rfoot{\ifthenelse{\value{page}=1}{\marginbox}{}}

% Margin
\newcommand{\marginbox}{%
	\vspace*{-240mm}\hspace*{105mm}\vbox{
		\vtop to 5mm {\vfil\usebox{\institutebox}\vfil}
		\vbox to 35mm {\vfil\usebox{\expeditorbox}\vfil}
		\vbox to 5mm {\vfil}
		\vbox to 5mm {\vfil\usebox{\datebox}\vfil}
		\vbox to 30mm {\vfil}
		\vbox to  35mm {\vfil\usebox{\researchcenterbox}\vfil}
		\vbox to 10mm {\vfil}
		\vbox to 20mm {\vfil\usebox{\leadingbox}\vfil}
		\vbox to 10mm {\vfil}
		\vbox to 30mm {\vfil\usebox{\bankbox}\vfil}
  }
}

% Institute
\newsavebox{\institutebox}
\newcommand{\institute}[2]{\savebox{\institutebox}{\parbox{100mm}{%
	\textbf{\scriptsize #1}}}}

% Expeditor
\newsavebox{\expeditorbox}
\newcommand{\expeditor}[5]{\savebox{\expeditorbox}{\parbox{100mm}{%
	{\scriptsize #1} \\{\scriptsize \textbf{#2 #3}}  \\{\scriptsize #4} \\[1ex]\scriptsize#5}}}

% Info of the research center
\newsavebox{\researchcenterbox}
\newcommand{\adress}[1]{\savebox{\researchcenterbox}{\parbox{100mm}{%
	\scriptsize\color{grey}\textbf{Helmholtz Zentrum München} \\German Research Center\\ for Environmental Health \\[2ex]#1 \\[2ex]+49 (0) 89 3187 (0) \\[4ex] info@helmholtz-muenchen.de \\www.helmholtz-muenchen.de}}}

% Info about leading of the research center
\newsavebox{\leadingbox}
\savebox{\leadingbox}{\parbox{100mm}{%
	\scriptsize\color{grey}\textsf{Aufsichtsratsvorsitzender: \\MinDir ‘in Bärbel Brumme-Bothe \\[4ex]Geschäftsführer: \\Prof. Dr. Günther Wess \\Dr. Nikolaus Blum \\Dr. Alfons Ehmsen}}}

% Date
\makeatletter
\newsavebox{\datebox}
\savebox{\datebox}{\parbox{100mm}{%
	\textsf{\footnotesize\@date}}}
\makeatother

% Bank
\newsavebox{\bankbox}
\savebox{\bankbox}{\parbox{100mm}{%
	\scriptsize\color{grey}\textsf{Registergericht: \\Amtsgericht München HRB 6466 \\USt-IdNr. DE 129521671 \\[4ex]Bankverbindung: \\Münchner Bank eG \\Konto-Nr. 2 158 620 \\BLZ 701 900 00 \\IBAN DE04 7019 0000 0002 1586 20 \\BIC GENODEF1M01}}}

% Recipient
\newcommand{\recipient}[1]{%
	#1}

\newcommand{\newpart}{\bigskip\bigskip\bigskip}
\newcommand{\newparagraph}{\bigskip}



\begin{document}
	% Set the name of the institute
	\institute{Institute for Epigenetics and Stem cells}

	% Set the expeditor (arguments are: title, firstname, lastname, position, additional infos like phone, email, webpage, etc.)
	\expeditor{Dr.}{Julien}{Pontabry}{Image analysis specialist}{%
		+49 (0) 89 3187 3398\\
		julien.pontabry@helmholtz-muenchen.de
	}

	% Address of the lab
	\adress{%
		Marchioninistraße 25\\
		D-81377 München\\
		Germany
	}

	% Recipient
	\recipient{%
		Dr. Julien Pontabry\\
		Institute for Epigenetics and Stem cells (IES)\\
		Marchioninistraße 25\\
		D-81377 München\\
		Germany
	}
	
	\newpart

	% Subject
	Subject: template for HMGU-styled letter with \LaTeX{}
	
	\newpart

	% Text
	Dear Madam, Mister,

	\newpart

	the \LaTeX{} version of the template for HMGU-styled letter is configurable with the following elements:
	\begin{itemize}
		\item \texttt{$\backslash{}$institute\{\textit{name}\}} identify the name of the institute (or department);
		\item \texttt{$\backslash{}$expeditor\{\textit{title}\}\{\textit{first-name}\}\{\textit{last-name}\}\{\dots\}} identify the author of the letter and is written in the margin;
		\item \texttt{$\backslash{}$adress\{\dots\}} contains the address of the lab;
		\item \texttt{$\backslash{}$recipient\{\dots\}} contains information about the recipient (name, compagny or lab, address, etc.);
		\item \texttt{$\backslash{}$newpart} creates a large vertical space to separate parts of the letter;
		\item \texttt{$\backslash{}$newparagraph} creates a small vertical space to separate paragraphs in the letter;
		\item you can change the date by using the standard \texttt{$\backslash{}$date} command (use \texttt{$\backslash$today} for current date).
	\end{itemize}

	\newparagraph

	This document needs the HMGU logo (file \textit{hmgu\_color\_englisch.eps}) and the following packages:
	\begin{itemize}
		\item \textit{fontenc} to create PDF with correct fonts (packages \textit{cm-super} or \textit{lmodern} are required to get a good result);
		\item \textit{graphicx} to display images (here logos);
		\item \textit{fancyhdr} to be able to modify headers and footers;
		\item \textit{color} to provide colored fonts;
		\item \textit{ifthen} to manage conditions;
		\item \textit{helvet} to write with Helvetica font;
		\item \textit{babel} with the english option to take into account the particularities of the language.
	\end{itemize}
	The file named \textit{hmgu-style.tex} is mandatory: it contains all the custom commands needed for compilation.

	\newparagraph

	Logos are included in EPS format. Therefore, the compilation has to be done with the following commands: \texttt{latex}, \texttt{dvitops} et \texttt{pstopdf}.

	\newparagraph

	This document corresponds to the version in date of \today. For all additional information, for bug report, or for suggestion, please send an email to julien.pontabry@helmholtz-muenchen.de.
	
	\newparagraph
	
	Your sincerely,

	\newpart

	Dr. Julien Pontabry
\end{document}
